% ============================================
% IMPLEMENTATION DETAILS
% ============================================
% 👤 PERSON 4: Implementation & Demo Experience
% 
% Instructions:
% - Technical stack (Python, libraries, GUI framework)
% - Key modules: simulator_full.py, integrate_l1_l2.py, GUI components
% - ML model training and deployment
% - Reference Figure 2 (GUI screenshots)
% - NO code dumps - explain what the code does, not how
% - Length: 0.7-1 page
% 
% NOTE: You own Figure 2 - create GUI screenshots
% ============================================

\section{Implementation Details}

[PERSON 4: Explain what users interact with and how it's built]

\subsection{Technology Stack}

ADRDS is implemented in Python 3.9+ with the following key technologies:
\begin{itemize}
    \item \textbf{Simulation Engine:} Custom behavioral simulator using multiprocessing 
          for parallel event generation
    \item \textbf{ML Framework:} scikit-learn for model training; supports Random Forest, 
          SVM, and neural network classifiers
    \item \textbf{GUI Framework:} [Tkinter/PyQt/Streamlit] for interactive visualization
    \item \textbf{Data Processing:} pandas and numpy for feature engineering
    \item \textbf{Visualization:} matplotlib/plotly for real-time charts
\end{itemize}

\subsection{Core Modules}

\textbf{Behavioral Simulator (\texttt{simulator\_full.py}):} Generates configurable 
ransomware behavioral sequences across attack stages. Users can adjust parameters 
such as encryption speed, file targeting strategies, and C2 communication frequency.

\textbf{Feature Integration (\texttt{integrate\_l1\_l2.py}):} Bridges L1 stage 
information with L2 behavioral events, producing feature vectors suitable for ML 
classification while maintaining interpretability.

\textbf{Demo Orchestrator (\texttt{demo\_complete.py}):} Coordinates simulation, 
detection, and visualization for seamless interactive demonstration experiences.

\subsection{Machine Learning Pipeline}

Models are trained on labeled datasets (see Section~\ref{sec:experiments}) using 
5-fold cross-validation. Feature selection employs mutual information scoring to 
identify the most discriminative behavioral indicators. Trained models are serialized 
for real-time inference during demonstrations.

\subsection{User Interface}

Figure~\ref{fig:gui} shows the main interface components:
\begin{itemize}
    \item Control panel for simulation parameter adjustment
    \item Real-time attack stage timeline
    \item Detection confidence meters with threshold sliders
    \item Alert notification panel
    \item Performance metrics dashboard
\end{itemize}

% PERSON 4: Insert Figure 2 here
% \begin{figure}[htbp]
% \centering
% \includegraphics[width=\columnwidth]{figures/gui_screenshot.png}
% \caption{ADRDS interactive GUI showing real-time detection and visualization.}
% \label{fig:gui}
% \end{figure}
